\section{Results}
\label{sec:results}

% Shell: structure and reporting format are fixed; numbers in \result{} come from
% `dockprepbench summarize/compare/figures`. Follow the benchmark-reporting conventions:
% denominators, both metrics, Wilson intervals, effect size before p-value, Holm-adjusted p.

\subsection{Baseline performance}

Of the 428 complexes, \result{N} were prepared and docked; \result{M} failed preparation,
\result{breakdown by cause}. With the baseline protocol, Vina placed the top-ranked pose
within \SI{2}{\angstrom} of the crystal ligand for \result{x\% [lo, hi]} of complexes, and
\result{y\% [lo, hi]} were also \pbvalid{}. \result{Comparison with the Vina results reported
by PoseBusters, listing protocol differences.}

\subsection{Run-to-run variation}

\todo{Baseline repeated with different seeds: spread of success across seeds. This sets the
  smallest effect we interpret.}

\subsection{Effect of each preparation choice}

\begin{figure}[t]
  \centering
  % \includegraphics[width=\textwidth]{success_rates.pdf}
  \fbox{\parbox[c][5cm][c]{0.9\textwidth}{\centering\result{Figure: success by arm with 95\%
    intervals (dockprepbench figures)}}}
  \caption{Top-1 and \pbvalid{} success for each experiment arm, with 95\% Wilson
    intervals. The dashed line marks the baseline top-1 success.}
  \label{fig:success}
\end{figure}

\result{One paragraph per arm, in the format: "Box padding of \SI{2}{\angstrom} changed top-1 success
from x\% [lo, hi] to y\% [lo, hi] ($\pm$d points; b vs c discordant complexes; Holm
$p = \dots$)."}

\begin{table}[h]
  \centering
  \caption{Paired comparison of each arm with the baseline. \result{Fill from
    \texttt{dockprepbench compare}.}}
  \label{tab:comparisons}
  \begin{tabular}{lrrrrr}
    \toprule
    Arm & Top-1 (\%) & $\Delta$ (pts) & Gained & Lost & Holm $p$ \\
    \midrule
    Baseline & \result{} & -- & -- & -- & -- \\
    \result{arms} & & & & & \\
    \bottomrule
  \end{tabular}
\end{table}

\subsection{How outcomes change: failure modes}

\result{Transition analysis: for each arm, the movements between failure classes behind the
net effect (e.g.\ scoring failures that became successes), with representative complexes.}

\subsection{Which complexes are sensitive to preparation}

\result{Relation between outcome changes and ligand properties (charge, flexibility, size)
and the presence of cofactors.}
